\section{结论}
\label{sec:conclusion}

本文在区分 Heretic 方向消融、上游 ARA 与本地适配的基础上，给出了校准式有限秩增量的目标、求解与规范化说明，并复核了 point-v1 的一次完整搜索。120 个候选全部完成，但无候选通过验收；验证集最低拒答关键词命中率为 54/100，对应首 token KL 为 0.089975。最低 KL 候选来自另一条记录，机器验收没有选中任何候选。独立 audit、重载分数及发布产物哈希均无记录，归档报告未导出适配器。这些事实表明，本次局部干预产生了指标变化，却没有满足预先指定的行为条件。

方法分析揭示了后续验证需要覆盖的环节：冻结局部输入与多层部署之间存在代理差距，关键词命中率与完整拒答语义并不相同，首 token 分布漂移也不足以衡量整体能力。trajectory-v2 已实现逐步对齐校准、部署增益、连续评分和隔离验收，但本归档没有它的效果数据。未来工作应先补齐同协议 Heretic 对照和多种子验证，再通过受控消融判断轨迹位置、目标形式和部署幅度各自的作用，并对锁定候选开展独立审计。本文不以这些待验证方向替代现有负结果，也不在本轮追加模型运行。
